\bookchapter{Bab 5: Membuat Ponsel Lebih Memahami Anda — Pengaturan Dasar}

\chapterintro{Luangkan beberapa menit untuk mengatur, peningkatan pengalaman sangat besar. Pengaturan ini seperti “menyetem” ponsel, membuatnya lebih sesuai dengan kebiasaan penggunaan Anda.}

\booksection{Suara dan Getaran}

Nada dering

\begin{itemize}
  \item \textbf{: Pilih yang bisa Anda dengar, tidak memalukan.}Mode getaran
  \item \textbf{: Saat senyap, ingatkan dengan getaran.}Mode jangan ganggu
  \item \textbf{: Bisa mengatur pengecualian, seperti mengizinkan panggilan dari kontak.}Pemisahan volume
  \item \textbf{: Volume media, volume nada dering, volume alarm bisa diatur terpisah.}2. Tampilan dan Kecerahan
\end{itemize}

\booksection{Tampilan dan Kecerahan}

\begin{itemize}
  \item \textbf{: Biarkan ponsel menyesuaikan secara otomatis berdasarkan cahaya lingkungan, praktis.}Mode perlindungan mata/night shift
  \item \textbf{: Menyaring cahaya biru, lebih nyaman saat melihat ponsel di malam hari.}Mode gelap
  \item \textbf{: Menghemat daya (sangat efektif untuk layar OLED), dan tidak silau di malam hari. Banyak aplikasi juga mendukung mode gelap.}Refresh rate layar
  \item \textbf{: Sebagian ponsel Android bisa diatur, semakin tinggi (seperti 120Hz) semakin halus scroll.}Kunci layar otomatis
  \item \textbf{: Atur 15 detik atau 30 detik, aman dan hemat daya.}Ukuran font
  \item \textbf{: Merasa huruf kecil? Perbesar sedikit.}3. Wallpaper, Tema, dan Layar Utama
\end{itemize}

\booksection{Wallpaper, Tema, dan Layar Utama}

Widget layar utama

\textbf{: Letakkan cuaca, jadwal, kontrol musik di layar utama, akses langsung. iOS dan Android sama-sama mendukung, tahan lama area kosong layar utama untuk menambahkan.}4. Pengaturan Efisiensi: Jadikan Ponsel “Otak Eksternal” Anda

\booksection{Pengaturan Efisiensi: Jadikan Ponsel “Otak Eksternal” Anda}

1. Mode Fokus / Jangan Ganggu: Rebut Kembali Perhatian yang Dirampas Notifikasi

\booksubsection{Mode Fokus / Jangan Ganggu: Rebut Kembali Perhatian yang Dirampas Notifikasi}

\begin{itemize}
  \item \textbf{: Bukan "kurangi melihat ponsel", tetapi "fokus saat harus fokus, hubungi lagi saat harus dihubungi".}Saran penggunaan
  \item \textbf{:}:\newline · Mode tidur: Aktif otomatis, hanya izinkan alarm;\newline · Olahraga / Membaca / Rapat: Sesuaikan dengan aktivitas saat ini.\newline Strategi lanjutan
  \item \textbf{: Gunakan otomatisasi + pemicu lokasi / waktu, biarkan mode beralih otomatis. Misalnya: iOS Shortcuts dapat mengatur "saat tiba di kantor otomatis masuk mode kerja".}2. Manajemen Notifikasi: Hentikan Titik Merah yang "Memanggil Tanpa Henti"
\end{itemize}

\booksubsection{Manajemen Notifikasi: Hentikan Titik Merah yang “Memanggil Tanpa Henti”}

\begin{itemize}
  \item \textbf{: 90\% push sebenarnya tidak Anda butuhkan.}Urutkan berdasarkan kepentingan
  \item \textbf{: Letakkan komunikasi interpersonal (telepon, SMS, pesan instan) di atas; berita, belanja, platform sosial bisa diatur sebagai "push senyap".}Notifikasi berkelompok
  \item \textbf{: Lipat notifikasi serupa bersama, hindari dibombardir puluhan pesan saat satu kali buka kunci.}3. Gerakan Cepat: Operasi Selangkah Lebih Cepat
\end{itemize}

\booksubsection{Gerakan Cepat: Operasi Selangkah Lebih Cepat}

\begin{itemize}
  \item \textbf{: Daya + volume bawah (iOS layar penuh)/Daya + volume bawah (Android umum).}Tangkapan layar panjang / tangkapan layar scroll
  \item \textbf{: Setelah tangkapan layar biasa klik opsi "tangkapan layar panjang".}Pindai kode sekali klik
  \item \textbf{: Tambahkan ikon "pindai" di Control Center / pengaturan cepat.}Ketuk belakang (iPhone)
  \item \textbf{: Di [Pengaturan] > [Aksesibilitas] > [Sentuh] > [Ketuk Belakang] ikatkan aksi seperti "tangkapan layar" atau "kunci layar".}Panggilan suara
  \item \textbf{: Siri / Xiao Ai Tongxue / Xiao Yi / Xiao Bu, sebut saja bisa cek cuaca, atur alarm.}4. Metode Input: Senjata Efisiensi Paling Diremehkan
\end{itemize}

\booksubsection{Metode Input: Senjata Efisiensi Paling Diremehkan}

\begin{itemize}
  \item \textbf{:}:\newline · iFlytek: Tingkat pengenalan input suara tertinggi, senjata untuk mencatat rapat;\newline · Gboard (Google): Sederhana, sinkronisasi lintas platform kuat, mendukung input geser;\newline · iOS bawaan: Privasi terbaik, integrasi mendalam dengan sistem.\newline Sakelar wajib
  \item \textbf{: Koreksi otomatis, kata saran, sinkronisasi kosakata cloud, input angka simbol dengan geser ke atas.}Catatan privasi
  \item \textbf{: Berhati-hati menggunakan fungsi "saran kalimat lengkap" atau "rekomendasi personalisasi", mungkin mengunggah konten input.}5. Layar Utama dan Multitugas: Lancarkan Alur Kerja
\end{itemize}

\booksubsection{Layar Utama dan Multitugas: Lancarkan Alur Kerja}

\begin{itemize}
  \item \textbf{: Balas pesan saat menonton video, catat saat browsing.}Jendela kecil / jendela mengambang
  \item \textbf{: WeChat / TikTok mendukung "mode jendela kecil", setelah balas langsung tutup.}Aplikasi yang sering digunakan disematkan
  \item \textbf{: Dock bar (bagian bawah layar) hanya letakkan 4 aplikasi paling sering, kurangi waktu pencarian.}Daftar Periksa "Hari Pertama Bekerja"
\end{itemize}

\begin{tcbinfo}{Saat mendapat ponsel baru atau pembaruan versi besar, disarankan luangkan 10 menit untuk meninjau pengaturan ini}
(1) Kata sandi kunci layar + biometrik (2) Notifikasi berkelompok + matikan push pemasaran (3) Ukuran font / refresh rate / mode gelap\newline (4) Metode input / browser / email default (5) Kontak darurat / kartu medis (6) Pengaturan cadangan\newline 5. Aksesibilitas: Kekuatan Teknologi untuk Kebaikan

\end{tcbinfo}

\booksection{Aksesibilitas: Kekuatan Teknologi untuk Kebaikan}

teknologi untuk kebaikan\textbf{: memungkinkan setiap orang menikmati kehidupan digital secara setara.}1. “Mata” Pengguna Low Vision — Panduan Lengkap Screen Reader

\booksubsection{“Mata” Pengguna Low Vision — Panduan Lengkap Screen Reader}

membacakan semua konten visual di layar\textbf{, pengguna mengoperasikan ponsel melalui gerakan dan suara. Berikut memperkenalkan penggunaan iOS “VoiceOver” dan Android TalkBack.}iOS VoiceOver

\begin{tcbinfo}{Jalur aktivasi}
\textbf{: [Pengaturan] > [Aksesibilitas] > [TalkBack], buka sakelar. Atau tahan tombol volume atas-bawah 3 detik untuk aktifkan cepat (sebagian merek mendukung).}Gerakan dasar

\textbf{:}:

\begin{enumerate}
  \item Usap dua jari: Scroll halaman penuh atas/bawah/kiri/kanan;
  \item Ketuk dua kali satu jari: Aktifkan elemen fokus saat ini (setara dengan "klik");
  \item Usap tiga jari: Beralih halaman;
  \item Ketuk tiga kali satu jari: Perbesar atau pilih (berbeda dengan ketuk tiga kali biasa).
  \item Pengaturan personalisasi
\end{enumerate}

\textbf{: [TalkBack] > [Pengaturan Text-to-Speech], bisa atur}kecepatan baca\textbf{(default sekitar 50\%, bisa dinaikkan ke 70-80\% untuk lebih lancar),}nada\textbf{, ganti}Android TalkBack

\end{tcbinfo}

\begin{tcbinfo}{Jalur aktivasi}
\textbf{: [Pengaturan] > [Aksesibilitas] > [TalkBack], buka sakelar. Atau tahan tombol volume atas-bawah 3 detik untuk aktifkan cepat (sebagian merek mendukung).}Gerakan dasar

\textbf{:}:

\begin{enumerate}
  \item Ketuk dua kali satu jari: Aktifkan elemen;
  \item Usap atas-bawah/kiri-kanan dua jari: Scroll;
  \item Ketuk sekali tiga jari: Masuk mode baca (baca berurutan);
  \item Ketuk sekali empat jari: Jeda/lanjutkan baca.
  \item Pengaturan personalisasi
\end{enumerate}

\textbf{: [TalkBack] > [Pengaturan Text-to-Speech], bisa atur}kecepatan\textbf{,},\textbf{, ganti}mesin suara\textbf{(disarankan pasang mesin Google Text-to-Speech, bacaan lebih alami).}2. Buat Huruf Kecil Besar, Buat Kabur Jelas — Penyesuaian Visual

\end{tcbinfo}

\booksubsection{Buat Huruf Kecil Besar, Buat Kabur Jelas — Penyesuaian Visual}

\booksubsection{Gerakan Perbesaran}

\begin{itemize}
  \item \textbf{: Pengguna biasa tidak dapat mengaktifkan mode pengembang lengkap, membutuhkan akun Apple Developer berbayar.}ketuk layar tiga jari\textbf{untuk memperbesar. Default zoom layar penuh, bisa pilih zoom jendela (jendela perbesaran lokal).}Cubit dua jari untuk atur perbesaran (1.5x hingga 15x)\textbf{, ketuk dua kali tiga jari untuk kunci rasio zoom saat ini.}Android
  \item \textbf{Android}ketuk layar tiga kali tiga jari\textbf{untuk aktifkan perbesaran sementara (maks 8 kali).}Teks Kontras Tinggi / Teks Tebal
\end{itemize}

\booksubsection{Teks Kontras Tinggi / Teks Tebal}

\begin{itemize}
  \item \textbf{: Pengguna biasa tidak dapat mengaktifkan mode pengembang lengkap, membutuhkan akun Apple Developer berbayar.}Android
  \item \textbf{Android}Filter Warna (Ramah Buta Warna)
\end{itemize}

\booksubsection{Filter Warna (Ramah Buta Warna)}

\begin{itemize}
  \item \textbf{: Pengguna biasa tidak dapat mengaktifkan mode pengembang lengkap, membutuhkan akun Apple Developer berbayar.}empat mode buta warna merah/hijau/biru/kuning\textbf{, serta opsi "Grayscale" dan "Inverse Warna". Disarankan lakukan tes penglihatan warna (bawaan sistem) sebelum memilih filter yang sesuai.}Android
  \item \textbf{Android}tiga mode buta warna\textbf{+ sakelar "Inverse Warna".}3. Deskripsi Gambar AI: Biarkan AI Menjadi "Mata" Anda
\end{itemize}

\booksubsection{Deskripsi Gambar AI: Biarkan AI Menjadi “Mata” Anda}

menjelaskan konten gambar dengan satu kalimat\textbf{, sangat membantu untuk mengenali CAPTCHA, membaca gambar, mengenali menu.}Dukungan Deskripsi Gambar AI Berbagai Sistem

\begin{tcbinfo}{iOS 15+}
\begin{itemize}
  \item \textbf{: Di [Pengaturan] > [Aksesibilitas] > [VoiceOver] > [Deskripsi Suara] aktifkan "Deskripsi Gambar". Sistem akan memanggil}model AI lokal\textbf{untuk menghasilkan deskripsi teks gambar otomatis (perlu iPhone 12 ke atas).}Aplikasi TalkBack versi baru
  \item \textbf{: Google TalkBack mengintegrasikan fungsi}tanya jawab gambar\textbf{di versi yang lebih baru, bisa bertanya "apa yang ada di gambar ini", AI menjawab real-time (tidak terlalu bergantung versi sistem Android, perbarui aplikasi TalkBack untuk mencoba).}Aplikasi pihak ketiga
  \item \textbf{: "Seeing AI" (iOS), "Lookout" (Android), "Xiao Yi Bang Bang" (Huawei) semua adalah asisten visual AI yang sangat baik, bisa mengenali barang, membaca dokumen, mengenali ekspresi wajah.}Skenario Penggunaan
\end{itemize}

\end{tcbinfo}

\begin{tcbpractice}{Coba bertanya pada foto street photography "berapa orang di gambar ini? apa yang mereka lakukan?", atau buka aplikasi ojek/pesan makanan, biarkan AI membaca nama menu di gambar menu —}
ini adalah “teknologi membuat hidup lebih setara” yang sebenarnya\textbf{.}.

\end{tcbpractice}

\booksubsection{“Telinga” Pengguna Tunarungu — Panduan Lengkap Caption Real-time}

fungsi speech-to-text tingkat sistem\textbf{, sangat berguna bagi pengguna tunarungu di lingkungan bising, perpustakaan, skenario bahasa asing.}Aktivasi dan Optimasi Caption Real-time

\begin{tcbinfo}{iOS}
\textbf{: Pengguna biasa tidak dapat mengaktifkan mode pengembang lengkap, membutuhkan akun Apple Developer berbayar.}Android 12+

\textbf{: Tekan}tombol volume\textbf{klik ikon “Caption Real-time” di popup; atau di [Pengaturan] > [Aksesibilitas] > [Caption Real-time] aktifkan.}Tips optimasi akurasi

\textbf{:}:

\begin{enumerate}
  \item lingkungan tenang\textbf{, jauhi AC, suara mesin;}Pengguna bicara
  \item artikulasi jelas, kecepatan sedang\textbf{memiliki tingkat pengenalan tertinggi;}Dialog multi-orang akan menurunkan akurasi, bisa gunakan mode "tampilkan satu orang saja";
  \item Mendukung bahasa Mandarin, Kanton, Inggris, dan bahasa mainstream lainnya,
  \item pengenalan bahasa minoritas relatif rendah\textbf{;}Saat bicara ada musik latar atau interferensi suara manusia, pengenalan akan terputus-putus.
  \item 5. Skenario Panggilan Dua Saluran: Speech-to-Text + Text-to-Speech
\end{enumerate}

\end{tcbinfo}

\booksubsection{Skenario Panggilan Dua Saluran: Speech-to-Text + Text-to-Speech}

tidak bisa mendengar juga tidak bisa bicara\textbf{(seperti kehilangan suara mendadak, masa pemulihan operasi tenggorok, stutter parah), ponsel masih bisa menjadi “juru bicara”.}Transkripsi Panggilan Real-time (TTY / RTT)

\begin{tcbinfo}{iOS}
\textbf{: Pengguna biasa tidak dapat mengaktifkan mode pengembang lengkap, membutuhkan akun Apple Developer berbayar.}teks yang Anda ketik juga akan dibacakan\textbf{ke lawan bicara.}Android

\textbf{Android}Kompatibilitas

\textbf{USB PD 3.1}:

\begin{itemize}
  \item \textbf{: Kedua pihak harus mendukung, teks diketik sambil tampil, tanpa jeda, operator mainstream telah mendukung;}TTY (Teletype)
  \item \textbf{: Memerlukan perangkat TTY khusus, telah secara bertahap digantikan RTT;}Fungsi
  \item speech-to-text\textbf{WeChat/QQ juga alternatif: tahan lama pesan suara "Ubah ke teks".}6. Fungsi "Kecil tapi Praktis" untuk Semua Orang
\end{itemize}

\end{tcbinfo}

\booksubsection{Fungsi “Kecil tapi Praktis” untuk Semua Orang}

\begin{itemize}
  \item \textbf{(iOS 14+): [Pengaturan] > [Aksesibilitas] > [Sentuh] > [Ketuk Belakang], bisa mengikat}ketuk dua/tiga kali\textbf{untuk memicu tangkapan layar, kunci layar, control center, dll. Sangat berguna saat kedua tangan membawa barang.}AssistiveTouch / bola mengambang
  \item \textbf{: iOS [Aksesibilitas] > [Sentuh] > [AssistiveTouch], Android [Aksesibilitas] > [Navigasi Mengambang]. Tombol virtual menggantikan gerakan atau tombol fisik, menyelamatkan saat tombol tidak berfungsi atau operasi satu tangan.}Kontrol Switch
  \item \textbf{: Mengontrol ponsel melalui sakelar eksternal, klik tunggal, gerakan kepala,}dirancang untuk pengguna disabilitas fisik\textbf{. iOS [Aksesibilitas] > [Kontrol Switch], Android [Aksesibilitas] > [Kontrol Sakelar].}Pengenalan suara
  \item \textbf{(iOS): [Aksesibilitas] > [Pengenalan Suara], setelah aktif ponsel akan "mendengar" bel pintu, tangisan bayi, alarm kebakaran, sirene, dll.,}saat dikenali akan getar + notifikasi layar\textbf{, sangat berharga bagi pengguna tunarungu dan saat tidur.}Notifikasi kedip LED
  \item \textbf{: iOS [Aksesibilitas] > [Audio/Visual] > [Kedip LED untuk Peringatan], Android [Pengaturan] > [Notifikasi] > [Notifikasi Lampu Kilat]. Saat panggilan masuk dan notifikasi lampu kilat berkedip,}bisa dirasakan tanpa suara di lingkungan bising\textbf{.}.
  \item \textbf{: [Pengaturan] > [Aksesibilitas] > [Audio] > [Audio Mono], menggabungkan saluran kiri-kanan. Pengguna dengan gangguan pendengaran satu telinga tidak akan kehilangan setengah informasi suara.}Filosofi Desain Aksesibilitas Apple
\end{itemize}

\begin{tcbinfo}{Apple berinvestasi besar di bidang aksesibilitas, setiap tahun memperbarui dan menambahkan fitur aksesibilitas. Situs web Apple memiliki halaman aksesibilitas khusus, detail penggunaan setiap fitur. Fitur aksesibilitas iOS biasanya lebih lengkap dan pengalaman lebih lancar daripada Android, ini juga alasan penting banyak pengguna disabilitas memilih Apple.}
Luangkan 30 Detik untuk Mencoba

\end{tcbinfo}

\begin{tcbpractice}{Cari waktu luang, lakukan pengalaman berikut}
Aktifkan

\begin{enumerate}
  \item VoiceOver / TalkBack\textbf{, coba tanpa mata menggunakan usap tiga jari untuk menjelajah WeChat;}Buka
  \item caption real-time\textbf{, bicara ke ponsel lihat akurasi pengenalan;}Gunakan
  \item ketuk tiga jari\textbf{untuk memperbesar teks hingga 200\%, rasakan keseharian pengguna low vision;}Matikan semua fitur sebelumnya, pikirkan "jika saya selalu membutuhkan fitur-fitur ini, seperti apa saya berharap didesain".
  \item Bab 6: Melindungi Kehidupan Digital Anda — Keamanan dan Privasi
\end{enumerate}

\end{tcbpractice}
